\subsection{Hallucination Audits and Source Grounding}
Dahl \emph{et al.} pose specific, verifiable questions about United States cases and compare model responses with known legal facts \cite{ref1}. Their taxonomy distinguishes several kinds of legal hallucination and their experiments show why fluent answers cannot be treated as reliable evidence. The contribution to this project is a motivation for independently checking identifiable claims. Its limitation is transfer: a United States audit neither establishes the error distribution in Indian filings nor measures this detector's performance. We consequently do not use its hallucination rates as target accuracy figures.

Magesh \emph{et al.} audit leading commercial AI legal research tools rather than general-purpose models \cite{ref2}. The study tests vendor reliability claims using legal research queries and evaluates generated responses. It motivates checking even systems that retrieve legal sources. However, these are end-to-end product results, while our controlled experiment concerns predefined mechanical changes to source-linked requests. Directly comparing those rates with our F1 score would compare different populations, labels, and tasks. We neither reproduce proprietary products nor claim superiority to their legal reasoning.

Lewis \emph{et al.} combine a pretrained generator with retrieval from nonparametric memory for knowledge-intensive tasks \cite{ref3}. This supports the architectural value of access to external evidence, but retrieval and verification remain distinct. The presence of a retrieved passage does not establish that the generated quotation reproduces it accurately. Gao \emph{et al.} introduce ALCE to evaluate answers with citations, including their correctness and completeness \cite{ref4}. Their work demonstrates why attaching a citation is weaker than supporting the full claim. General question-answering corpora do not provide the Indian court, reporter-alias, and numbered-paragraph annotations needed here. Our implementation therefore avoids claiming that a text lookup solves citation entailment.

\subsection{Fact Verification and Evaluation Limits}
FEVER introduces a large claim-verification dataset using Wikipedia evidence and supported, refuted, or insufficient-information labels \cite{ref5}. Its distinction between a negative verdict and insufficient evidence informs this project's unresolved state. Nevertheless, evidence selection in encyclopedic text is not a substitute for identifying the deciding court's voice or the scope of a legal holding. SciFact pairs expert-written scientific claims with abstracts and supporting or refuting evidence \cite{ref6}. The expert-annotation component is particularly relevant: automatically generating legal error labels is convenient, but cannot replace independent legal judgment. Neither benchmark is used as a surrogate Indian evaluation set.

FactCC studies factual consistency between source documents and abstractive summaries using synthetic transformations \cite{ref7}. This supplies a useful precedent for controlled perturbations while exposing their limitation: good performance on artificial transformations does not ensure detection of naturally occurring errors. SelfCheckGPT instead uses consistency across sampled model generations to detect hallucinations without an external reference collection \cite{ref8}. Such sampling provides a different reliability signal from this project's explicit judgment evidence. When an original source is available, directly inspecting the claimed words offers a simpler, explainable check; it still leaves the proposition's legal meaning unresolved.

TRUE standardizes evaluation of factual-consistency metrics across annotated datasets \cite{ref19}. Its emphasis on example-level assessment and variation between tasks supports preserving per-case predictions, confusion counts, and category breakdowns. A single aggregate number can conceal different failure costs. In legal checking, a false accusation that a case is invented and an unnecessary human-review warning are not interchangeable outcomes. We therefore report triage and four-state agreement separately rather than relabelling every warning as a proven false citation.

\subsection{Legal Language Resources}
LegalBench organizes diverse legal reasoning tasks through a collaborative benchmark \cite{ref9}. It supports task-specific capability evaluation, but a broad legal benchmark does not supply an operational Indian citation-indexing pipeline. LexGLUE standardizes multiple English legal language-understanding datasets and compares generic with legal-oriented models \cite{ref10}. Its findings support domain-sensitive evaluation. Document classification performance, however, says little about exact quotation fidelity, a missing paragraph identifier, or an authority's subsequent treatment.

Malik \emph{et al.} introduce ILDC for court judgment prediction and explanation using Indian Supreme Court documents \cite{ref11}. The source domain is closely relevant, and the work demonstrates that Indian judgments can support reproducible NLP research. Its target remains different: predicting an outcome from a judgment is not verifying a quotation supplied by a reader. This project therefore uses an attributed public source and defines a separate integrity task instead of presenting an outcome-prediction score as evidence of citation accuracy.

BERT learns contextual language representations through bidirectional pretraining \cite{ref12}. LEGAL-BERT investigates pretraining and adaptation for legal text \cite{ref14}, showing why domain mismatch deserves attention. Neither pretraining objective certifies legal support. A legal-domain encoder could improve candidate retrieval, but its score would still require task-specific evaluation and would not replace source-level metadata checks. We use deterministic checks for the mechanical requirements and reserve pretrained representations for a separately reported retrieval experiment.

\subsection{Retrieval Algorithms and the Remaining Gap}
Sentence-BERT adapts transformer representations for efficient sentence-level similarity search \cite{ref13}. The Transformer architecture supplies the underlying contextual attention mechanism \cite{ref18}. These approaches make vector comparison practical, but generic sentence similarity can reward paraphrases that differ in a legally decisive detail. Dense Passage Retrieval learns question--passage representations for open-domain question answering \cite{ref16}; its relevance criterion and candidate collection differ from within-judgment quotation localization. The current experiment therefore compares candidate ranks without treating dense similarity as an integrity verdict.

Robertson and Zaragoza explain probabilistic relevance ranking, including BM25's term-frequency saturation and document-length normalization \cite{ref15}. BM25 is attractive here because a quoted fragment usually shares vocabulary with its source and the score is inexpensive to compute. BEIR evaluates retrieval methods across heterogeneous zero-shot tasks and shows that effectiveness depends on domain and dataset \cite{ref17}. Together these results motivate a measured lexical-versus-dense comparison rather than assuming that GPU use necessarily improves retrieval.

The gap addressed is an inspectable workflow connecting Indian citation aliases, original paragraph numbers, quotation fidelity, mechanical metadata checks, and uncertainty reporting. We do not claim this is the first or only such system. The literature instead supports a bounded engineering conclusion: evidence presentation, explicit task definitions, and abstention remain necessary even with strong language models. The project contributes a runnable implementation and transparent measurement of that subset. Independently annotated natural errors, voice attribution, legal entailment, and subsequent-treatment analysis remain separate research problems.
